\section{ИССЛЕДОВАНИЕ И ОЦЕНКА КАЧЕСТВА РАЗРАБОТАННОЙ СИСТЕМЫ}

\subsection{Цели и задачи испытаний}

После реализации программной системы необходима проверка того, что полученное
решение соответствует сформулированным требованиям и сохраняет корректность
работы в ключевых для предметной области сценариях. Для системы индексации
ончейн-данных такая проверка должна охватывать не только общие API-функции, но
и устойчивость к reorg, корректность работы с mempool, согласованность
хранилища и работоспособность управляющих механизмов.

Целью испытаний является подтверждение пригодности разработанной системы
\texttt{Bitcoin Blockchain Indexer} к решению задач индексации, хранения и
выдачи ончейн-данных. Для достижения этой цели
необходимо решить следующие задачи:
\begin{enumerate}
    \item проверить корректность функций управления заданиями индексирования и взаимодействия с API;
    \item проверить согласованность записи подтвержденных данных в хранилище;
    \item оценить корректность обработки mempool и реорганизаций цепи;
    \item подтвердить работоспособность механизмов мониторинга и диагностики;
    \item сопоставить фактические результаты с установленными критериями проверки.
\end{enumerate}

\subsection{Программа и методика испытаний}

Испытания системы строятся на сочетании интеграционного тестирования,
функциональной проверки API и анализа критериев приемки. В качестве
основного подхода используются тесты, запускаемые в окружении с контейнером
PostgreSQL через \texttt{testcontainers}, что позволяет приблизить условия
проверки к реальной эксплуатации, сохраняя воспроизводимость тестовых
сценариев \cite{testcontainers_docs}.

В рамках программы испытаний проверяются следующие группы сценариев:
\begin{enumerate}
    \item сценарии жизненного цикла заданий индексирования: создание, запуск, пауза,
    возобновление, остановка и обработка ошибочных переходов;
    \item сценарии API nodes и data API, включая запросы балансов, UTXO,
    транзакций, блоков и mempool-данных;
    \item сценарии работы контура индексации и проверки идемпотентности
    записи подтвержденных блоков;
    \item сценарии фоновых процессов обработки, включая синхронизацию mempool и обработку reorg;
    \item сценарии наблюдаемости, связанные с публикацией метрик и
    диагностической информации о состоянии узлов и заданий индексирования.
\end{enumerate}

Методически испытания ориентированы на подтверждение ключевых свойств системы:
функциональной корректности, согласованности данных, восстанавливаемости после
сбоев и эксплуатационной наблюдаемости. Для итоговой оценки используется
перечень критериев приемки, фиксирующий степень выполнения отдельных
требований.

\subsection{Проверка корректности функциональных возможностей}

Функциональная корректность системы подтверждается набором интеграционных
тестов и реализованными прикладными сценариями взаимодействия с API.
Проверяются операции управления заданиями индексирования, получение списка узлов, запросы
подтвержденных транзакций, балансов адресов, UTXO и блоков, а также
пограничные ситуации, связанные с ошибками авторизации, отсутствием объекта
или нарушением ограничений переходов состояний.

С точки зрения предметной области особенно важны результаты проверки следующих
свойств:
\begin{enumerate}
    \item задания индексирования обладают формализованным жизненным циклом и корректно меняют
    состояние через API;
    \item подтвержденные данные сохраняются в БД в виде связанных сущностей:
    блоков, транзакций, входов, выходов и агрегатов адресного уровня;
    \item mempool-данные доступны как отдельный класс сущностей и не смешиваются
    с представлениями подтвержденной цепи;
    \item выдача данных через API ориентирована на прикладные выборки, а не на
    сырые RPC-ответы.
\end{enumerate}

Таким образом, функциональное тестирование показывает, что разработанная
система реализует основной набор возможностей, предусмотренный
сформулированными требованиями к программному комплексу, и обеспечивает
прикладную форму доступа к данным, отличную от прямой работы с RPC узла.

\subsection{Оценка устойчивости к reorg и изменениям состояния mempool}

Для систем индексации ончейн-данных устойчивость к изменениям состояния сети
является одним из ключевых критериев качества. В рассматриваемом проекте это
свойство проверяется через сценарии, связанные с reorg и динамикой mempool.

Испытания reorg подтверждают, что при обнаружении расхождения между сохраненной
в БД цепью и актуальными данными узла система:
\begin{enumerate}
    \item помечает затронутые блоки как неканонические;
    \item переводит связанные транзакции в статус \texttt{orphaned};
    \item пересчитывает производные агрегаты по оставшейся канонической ветке;
    \item откатывает прогресс задания индексирования до согласованной высоты.
\end{enumerate}

Испытания mempool подтверждают, что фоновый процесс синхронизации корректно обрабатывает появление
новых неподтвержденных транзакций и помечает исчезнувшие как
\texttt{dropped}. При этом текущие балансы подтвержденной цепи и подтвержденные \gls{utxo}
не искажаются за счет временного состояния mempool. Такое разделение
представлений принципиально важно для достоверности прикладных ответов.

Следовательно, система демонстрирует устойчивость к двум наиболее характерным
для предметной области видам изменчивости: локальному пересмотру канонической
цепи и колебаниям состава неподтвержденных транзакций.

\subsection{Оценка эксплуатационных характеристик и наблюдаемости}

Помимо функциональной корректности, для практического использования системы
важны ее эксплуатационные свойства. В разработанном решении они обеспечиваются
через контейнеризованное развертывание, конфигурирование из внешних источников,
JSON-логирование и публикацию метрик Prometheus \cite{prometheus_docs}.

К числу значимых эксплуатационных характеристик относятся:
\begin{enumerate}
    \item воспроизводимость запуска серверного приложения и базы данных в контейнерном окружении;
    \item возможность мониторинга прогресса заданий индексирования и текущей высоты цепи;
    \item наблюдение за RPC-запросами, задержками и ошибками записи в БД;
    \item наличие диагностической информации по состоянию узлов;
    \item пригодность системы к сопровождению через API, CLI и административную панель.
\end{enumerate}

В проекте экспортируются метрики текущей высоты цепи, прогресса заданий индексирования, лага
индексации, количества обработанных блоков и транзакций, а также статистика
RPC-запросов и ошибок. Это позволяет использовать систему как основу для последующей эксплуатационной
интеграции.

\subsection{Анализ результатов испытаний}

Анализ выполненных испытаний показывает, что основная часть требований к системе
реализована и подтверждена кодом, интеграционными тестами и критериями
приемки. Наиболее полно подтверждены работоспособность схемы хранения,
жизненный цикл заданий индексирования, корректность записи подтвержденных данных, обработка reorg,
разделение mempool и подтвержденной цепи, а также наличие REST API, CLI и
административной панели.

Вместе с тем анализ результатов испытаний показывает наличие направлений,
требующих дальнейшего расширения экспериментальной проверки. К ним относятся
подтверждение защищенного эксплуатационного контура взаимодействия,
полномасштабные end-to-end испытания с реальным Bitcoin Core/regtest и более
детальная фиксация итоговых показателей покрытия серверной части тестами.
Указанные вопросы не ставят под сомнение достигнутую работоспособность
системы, однако определяют перспективы углубления ее верификации.

В целом результаты испытаний позволяют сделать вывод о том, что разработанная
система успешно решает основную задачу индексации ончейн-данных, а выбранные
архитектурные и проектные решения являются обоснованными с точки зрения
предметной области и дальнейшего развития проекта.

\section{ПРАКТИЧЕСКАЯ ЗНАЧИМОСТЬ И ПЕРСПЕКТИВЫ РАЗВИТИЯ СИСТЕМЫ}

\subsection{Практическая значимость разработанного решения}

Практическая значимость разработанной системы определяется тем, что она
формирует специализированный слой доступа к ончейн-данным Bitcoin, пригодный
для использования прикладными и аналитическими сервисами. В отличие от прямой
работы с узлом сети, такое решение предоставляет уже нормализованные и
подготовленные данные, что снижает сложность интеграции и уменьшает нагрузку
на инфраструктуру источника данных.

Для практического применения наиболее важны следующие свойства разработанной
системы:
\begin{enumerate}
    \item наличие прикладного API для доступа к блокам, транзакциям, UTXO,
    балансам адресов и состоянию заданий индексирования;
    \item возможность управляемой индексации с контролем прогресса и ошибок;
    \item наличие административных и вспомогательных средств сопровождения;
    \item готовность к контейнерному развертыванию и мониторингу;
    \item сохранение согласованности данных подтвержденной цепи при reorg и при работе с mempool.
\end{enumerate}

\subsection{Возможности применения системы в прикладных сервисах}

Разработанная система может использоваться в составе различных программных
решений, работающих с блокчейн-данными Bitcoin. В частности, она подходит для:
\begin{enumerate}
    \item аналитических сервисов, выполняющих выборки по адресам, транзакциям и блокам;
    \item административных и мониторинговых панелей, отображающих состояние узлов и заданий индексирования;
    \item внутренних корпоративных сервисов, которым требуется контролируемый
    доступ к ончейн-данным без прямого обращения к RPC узла;
    \item исследовательских и учебных стендов, демонстрирующих процессы
    индексации, хранения и предоставления блокчейн-данных.
\end{enumerate}

Наличие CLI и административной панели расширяет область практического
использования системы, поскольку позволяет применять ее как в автоматизированных
сценариях, так и в задачах оперативного сопровождения.

\subsection{Ограничения текущей версии системы}

Несмотря на достижение основной цели работы, текущая версия системы
имеет ряд ограничений. К ним относятся:
\begin{enumerate}
    \item неполное подтверждение end-to-end сценариев с реальным Bitcoin Core;
    \item необходимость дальнейшего подтверждения защищенного эксплуатационного контура взаимодействия;
    \item ограниченная глубина экспериментальной оценки производительности;
    \item применение общей стратегии индексации без специализированной
    оптимизации для отдельных пользовательских режимов, например \texttt{address\_list};
    \item отсутствие расширенного эксплуатационного alerting и более детальных
    сценариев деградации узлов.
\end{enumerate}

Указанные ограничения не снижают значимости достигнутого результата, но
определяют границы текущего этапа разработки и подтверждают, что система
находится на стадии, ориентированной на дальнейшее инженерное развитие.

\subsection{Направления дальнейшего развития системы}

Перспективы развития разработанной системы связаны как с углублением
функциональности, так и с расширением эксплуатационных возможностей. Наиболее
естественными направлениями дальнейшей работы являются:
\begin{enumerate}
    \item проведение полноценных end-to-end испытаний с реальным узлом Bitcoin Core;
    \item развитие защищенного контура взаимодействия серверной части и внешних клиентов;
    \item оптимизация производительности индексирования и прикладных выборок;
    \item расширение режимов работы заданий индексирования и стратегий фильтрации по адресам;
    \item развитие средств мониторинга, диагностики и автоматического реагирования;
    \item подготовка архитектурной базы для дальнейшего масштабирования решения.
\end{enumerate}

С учетом выбранного архитектурного подхода указанные направления могут быть
реализованы поэтапно без отказа от уже созданной модели хранения и структуры
основных модулей.

\subsection{Выводы по результатам работы}

Разработанная система индексации ончейн-данных обладает практической ценностью
как самостоятельный программный компонент, обеспечивающий контролируемое
получение, хранение и выдачу блокчейн-данных Bitcoin. Ее применение позволяет
снизить сложность интеграции прикладных сервисов с блокчейн-инфраструктурой и
создает основу для дальнейшего развития в сторону более полного промышленного
решения.

\conclusion

В выпускной квалификационной работе решена задача проектирования и разработки
системы индексации ончейн-данных сети Bitcoin. В процессе выполнения работы
были последовательно рассмотрены теоретические основы предметной области,
проанализированы особенности блокчейна Bitcoin как источника данных,
сформированы требования к системе, спроектированы архитектура программного
комплекса и модель хранения, а также описана реализация ключевых модулей
индексации, управления, мониторинга и прикладной выдачи данных.

В результате выполненной работы создан программный комплекс
\texttt{Bitcoin Blockchain Indexer}, реализующий индексацию блоков и
транзакций, поддержку текущего состояния \gls{utxo}, адресных балансов,
обработку mempool и reorg, управление заданиями индексирования, а также предоставление
данных через REST API, административную панель и CLI. Проведенные испытания
подтвердили работоспособность основных функций системы и показали обоснованность
выбранных архитектурных и проектных решений.

Практическая значимость работы состоит в создании основы для дальнейшего
развития инфраструктурного решения, предназначенного для использования в
прикладных и аналитических сервисах, работающих с ончейн-данными Bitcoin.
Полученные результаты могут быть использованы как для дальнейшей инженерной
доработки проекта, так и для развития исследований в области системной
индексации блокчейн-данных.
